% !PRINTED 403
\weberpage{403}
%<<CONT
\[
\frac{1}{1}, \frac{2}{1}, \frac{7}{4}, \frac{16}{9}, \frac{23}{13}, \frac{2201}{1244}, \frac{4425}{2501} \cdots
\]
If we convert the two last fractions into decimal fractions, we obtain for $x$ the two limits:
\[
1{,}7692926; \quad 1{,}76929228,
\]
and the error in the latter, a number slightly too small, is less than
\[
0{,}00000016.
\]
That one here obtains a good result after few steps rests on the circumstance that a larger partial denominator 95 occurs fairly early. In most cases one obtains small numbers for the partial denominators; the denominators of the convergents then grow slowly, and one only laboriously obtains a result of greater accuracy.

The finding of larger partial denominators, e.g. 95, is made easier by the remark that, when a root of an equation has a large value, the negative coefficient of the second-highest power of the unknown gives an approximation to this root, even if only a rough one, as e.g. in the next to last of the above equations 94 to 95.

\begin{center}
\S.\,130.\\[0.3em]
\gesperrt{Rational Roots of Equations with Integral Coefficients.}\\
\gesperrt{Reducible Equations.}
\end{center}

We close this section with a consideration of equations with integral coefficients which, even though it does not stand in immediate relation to the theory of continued fractions, yet most fittingly finds a place here.

When one has to do with equations whose coefficients are rational numbers, one will first look for possible rational roots. If the equation
\[
f(x) = a_0 x^n + a_1 x^{n-1} + a_2 x^{n-2} + \cdots + a_{n-1} x + a_n = 0
\]
has rational coefficients, we can always assume that these coefficients are integers; to reduce the case of fractional coefficients to this one, one need only multiply the whole equation by a common multiple of all the denominators.
